\documentclass[11pt,a4paper]{article}
\usepackage[utf8]{inputenc}
\usepackage{graphicx}
\usepackage[left=2.5cm,top=3cm,right=2.5cm,bottom=3cm,bindingoffset=0.5cm]{geometry}
\usepackage{AEDLogica, AEDEspecificacion, AEDTADs}
\usepackage{caratula}
\newlength{\miSangria}
\setlength{\miSangria}{2em}

% Asi pueden escribir nuevos comandos. 
% Este por ejemplo asegura q los nombres 
% que figuren con una tipografia diferenciada  
\newcommand{\Tipo}[1]{\mathsf{#1}}
% la sintaxis es \newcommand{\nombreDeLaMacro}[cantidadDeParametros]{Lo que va ser remplazado por el macro} 
\newcommand{\norm}[1]{\vert#1\vert}


\begin{document}

\section{Titulo del Proc ACA}
\vspace{1.0cm}
\noindent Transacciones \textbf{ES} Tuple $\tupla{
		UsuarioOrigen: \Tipo{\Z}, \\
		UsuarioDestino: \Tipo{\Z},
		TipoDeTransaccion: \Tipo{\Z},
		Millas: \Tipo{\R}, \\
		NombrePromo: \Tipo{\str},
		FactorPromo: \Tipo{\R}
	}$

\noindent \comentario{Tipo de transacciones es algo que yo pongo a mano creo, ejemplo 0:Transaccion por referidor, 1: Transferencia, 2: Canjeo de puntos}


\begin{tad}{Airmiles}



	\obs{Usuarios}{\Diccionario{id}{\tupla{MillasTotales: \Tipo{\Z}, MillasDisponiles: \Tipo{\Z}}}}

	\noindent \obs{HistorialDeTransacciones}{\seq{Transacciones}}
	\\
	\texttt{obs} Promociones : dict \{ \\
	\hspace*{1.5em} NombrePromocion: $\str$, \\
	\hspace*{1.5em} $\seq{\Tipo{FactorDePromocion},\Tipo{topeDePromocion}, \Tipo{millasRestantes}}$ \\
	\}


	\begin{proc}{paresTransferidoresExclusivos}
		{\In Am: Airmiles}
		{\seq{\tupla{\Tipo{\Z}, \Tipo{\Z}}}}
		\requiere{True}

		\aseguraLargo{
            NoHayRepetidos(res) \land
            \\
            NoHayPermutacion(res) \land
            \\
            TieneReciprocoHistorial(AM.Historial,res) \land
            \\
            SonParejaExclusiva(AM.Historial,res)
        }
	\end{proc}

    \predLargo{NoHayRepetidos}{res}{
        \paraTodo{i}{\Tipo{\Z}}{0 \le i < \norm{res} \implicaLuego \paraTodo{j}{\Tipo{\Z}}{0 \le j < \norm{res} \yLuego j \neq i} \implicaLuego (res[i]_{0} \neq res[j]_0) \land res[i]_{1} \neq res[j]_{1}}
    }

    \comentario{Cual es la forma correcta de hacer esto? algo asi pero tengo que indicar el orden en el que estoy poniendo las cosas?}
    \\
    \predLargo{NoHayPermutacion}{res}{
        \paraTodo{t}{\Tipo{\tupla{\Tipo{\Z},\Tipo{\Z}}}}{t \in res \implicaLuego \tupla{{t_1},t_{0}} \notin res}
    }

    \vspace{0.5cm}
    \predLargo{TieneReciprocoHistorial}{Historial,res}
    {\paraTodo{x}{\tupla{\Tipo{\Z},\Tipo{\Z}}}{x \in res \implica HayReciproco(historial)}}

    \vspace{0.5cm}
    \predLargo{HayReciproco}{Historial,t}
    {
        \existe{i}{\Tipo{\Z}}{0 \le i < \norm{historial} \yLuego 
        \\
        (Historial[i].usuarioOrigen = t_{1} \land Historial[i].usuarioDestino = t_{0})}
    }

    \vspace{0.5cm}
    \noindent \comentario{Revisar}
    \\
    \predLargo{SonParejaExclusiva}{Historial,res}
    {
        \paraTodo{t}{\tupla{t_{0}: \Tipo{\Z},t_{1} :\Tipo{\Z}}}{ t \in res \implicaLuego
        \neg (hayOtroUsuario(t0,t1,Historial) \lor 
        \\
        hayOtroUsuario(t1,t0,Historial))
        }
    }

    \vspace{0.5cm}
    \predLargo{HayOtroUsuario}{t0,t1,Historial}
    {
        \paraTodo{i}{\Tipo{\Z}}{ 0 \le i < \norm{Historial} \implicaLuego (Historial[i].UsuarioOrigen = t0 \land
        \\ 
        Historial[i].UsuarioDestino \neq t1)}
    }

\end{tad}
\end{document}